\documentclass[10pt,a4paper]{amsart}
\usepackage[margin=19mm]{geometry}
\usepackage{amsmath,amssymb,mathtools}
\usepackage[scr=rsfs,cal=euler]{mathalfa}
\usepackage{xfrac}
\usepackage[unicode,hidelinks,bookmarks=false]{hyperref}
\newcommand{\ie}{i.e.\ }
\newcommand{\cf}{cf.\ }
\newcommand{\resp}{respectively\ }
\newcommand{\Subsection}{\subsection}
\providecommand{\nobreakdash}{\nobreak\hbox{-}}
\setlength{\emergencystretch}{2em}
\multlinegap0pt
\usepackage{bm}
\usepackage{bbm}
\usepackage{dsfont}
\usepackage{xspace}
\usepackage{cancel}
\usepackage{mathtools}
\usepackage{amsthm}
\makeatletter
\@ifundefined{prop}{\newtheorem{prop}{Proposition}}{}
\@ifundefined{theorem}{\newtheorem{theorem}{Theorem}}{}
\makeatother
\title[On the word problem for just infinite groups]{On the word problem for just infinite groups}
\author{Alexey Talambutsa}
\date{Source: arXiv:2512.24266v3 (CC BY 4.0)}
\begin{document}
\maketitle

\noindent\emph{Source and licence.} On the word problem for just infinite groups. Version arXiv:2512.24266v3 (CC BY 4.0), distributed under the Creative Commons Attribution 4.0 International licence, as recorded in the arXiv licence metadata for this exact version and shown on the abstract page https://arxiv.org/abs/2512.24266v3. Full rights notice: licenses/arxiv-2512.24266-CC-BY-4.0.txt.\par\smallskip

\noindent\textit{Editorial scope.} Counted unit: the Theorem whose source label is \texttt{th-undecidable} in the article (source0.tex lines 323--326, source label \texttt{th-undecidable}), delivered together with the article's own complete original proof of it (source0.tex lines 327--380, one \texttt{proof} environment of 54 lines). No mathematics is written, repaired or re-derived: statement and proof are the article's own text line for line. Required context retained uncounted: it-case (lines 261--268); prop-shaded-set (lines 275--278). Every definition, display and label of those lines is retained unchanged. Omitted from this package and not redistributed: 1--260, 269--274, 279--322, 381--479. Nothing inside a retained excerpt is omitted or altered: every retained body is delivered exactly as the article's lines are written, so no omission receipt of any kind accompanies this package. Citations: the retained excerpts carry no citation command at all, so no bibliography entry of the article is transcribed and no reference is redistributed. Reference adaptation: none was needed and none was made; every reference inside the delivered unit resolves inside it, so no unresolved reference or citation is produced. Printed slips kept literal: no printed word, symbol, hypothesis or number of the counted statement or of its proof was altered. Layout and class: the article's own class is \texttt{amsart} with packages including babel, inputenc, fontenc, textcase, pdflscape, amsmath, amssymb, latexsym, graphics, epsfig, color, hyperref, geometry; the wrapper is the lane's pinned \texttt{amsart} editorial wrapper at 10pt on A4, which restates the theorem-like environments the retained text needs and adds the article's own macro definitions that the retained text needs (none), with extra wrapper packages (bm, bbm, dsfont, xspace, cancel, mathtools). The delivered numbering is the unit's own. Assets: no retained span contains a figure environment, a graphics-inclusion command, a PDF-page inclusion, a table or a TikZ picture; no image file is copied into this package, the package declares no asset dependency and no image file is redistributed with it. Licence: version arXiv:2512.24266v3 of this article is distributed under the Creative Commons Attribution 4.0 International licence, as recorded in the arXiv licence metadata for this exact version and shown on the abstract page https://arxiv.org/abs/2512.24266v3. A full rights notice is delivered at licenses/arxiv-2512.24266-CC-BY-4.0.txt. Characters: every retained excerpt is pure ASCII, so the delivered engine, XeLaTeX with its default Latin Modern text font, reports no missing character.
\begin{equation}
\label{it-case}
I_t =
 \begin{cases}
   [l_t,r_t] &\text{if $l_t>\max \{r_i \mid i<t\}$},\\
   [l_t,\max \{r_i \mid i\leq t\}] &\text{otherwise}.
 \end{cases}
\end{equation}

\begin{prop}
\label{prop-shaded-set}
There exists a recursively enumerable sequence of segments $\alpha=\{[l_t,r_t]\}_{t=0}^\infty$ for which the shaded set $\mathcal I(\alpha)$ is recursively enumerable and not recursive.
\end{prop}

\begin{theorem}
For each group $W_C$ there exists a countably generated recursively enumerable presentation $\langle S \mid \mathcal U \rangle$ with an undecidable word problem.
\label{th-undecidable}
\end{theorem}

\begin{proof}
First, we describe the standard way to construct for the group $W_C$ a recursively enumerable presentation $\langle S \mid \mathcal R \rangle$ with decidable word problem. Then we will show how to modify this process using the recursively enumerable sequence $\alpha$ from Proposition~\ref{prop-shaded-set} to obtain the desired group presentation $\langle S \mid \mathcal U \rangle$ with undecidable word problem.

\smallskip

Let us choose a finite presentation $C=\langle a_1,\ldots,a_d \mid q_1=1, \ldots,q_m=1 \rangle$ for the initial perfect finite permutational group $C$ acting on the set of size $r$. 

Due to the associativity of permutational wreath products, one can as well write the group $W_C$ as $\lim_{n\to \infty} G_n$, where $G_0=\{1\}$ and $G_{n+1}= C \wr G_{n}$. Thus, $G_{n+1}$ is a semidirect product $B_n \rtimes G_n$, where the normal subgroup $B_n$ is a direct product of $|q^n|$ copies of the group $C$. The permutational group $G_n$ acts on $B_n$ by permuting these copies. 

To obtain a presentation for the group $G_{n+1}$ we take a presentation of $G_n$ and extend it to a presentation of the free product $G_{n}*B_n$ by adding $dq^n$ new generators and relations in them: the $mq^n$ defining relations for $q^n$ copies of the group $C$ and commutations between the generators of different copies. After that we add finitely many relations that define the action of generators for $G_n$ on the copies of $C$ and thus obtain $G_{n+1}$ as a quotient of $G_{n}*B_n$. In this construction the canonical projecting homomorphism $\pi_{n+1}: G_{n+1}=B_n \rtimes G_n \to G_n$ is realized by sending the $dq^n$ new generators to $1$. 

If we define the recursive function
\begin{equation}
f(0)=0, \ f(1)=d \quad \text{and} \quad  f(n+1)=f(n)+dq^n,
\end{equation}
then the process above gives a presentation of the group $G_n$ with $f(n)$ generators, and we will denote it as $GP_n=\langle s_1,s_2,\ldots,s_{f(n)} \mid R_n \rangle$. Note that the presentation $GP_n$ can be considered as the initial part of the presentation $GP_{n+1}$, which also agrees with the canonical embedding $\psi_n:G_n\to G_{n+1}$. Passing to the limit as $n\to\infty$, one obtains a countable set of generators and a countable set of relations that define the limit presentation $W_C=\langle S \mid \mathcal R \rangle$.


%The description of the presentation $\langle S \mid \mathcal R \rangle$ above shows that $G_n\cong\langle s_1,\ldots, s_{E(n)} \mid R_n \rangle$, where $R_n$ are some first relations of $\mathcal R$ that are written in the alphabet $\{ s_1,\ldots, s_{E(n)}\}^\pm$.

%\smallskip

%Consider $k,n\in \mathbb N$ such that $k<n$. A presentation of the group $G_k$ can be obtained from the presentation $GP_n=\langle s_1,\ldots, s_{f(n)} \mid R_n \rangle$ by adding relations $s_{E(k)+1}=1,\ldots,s_{E(n)}=1$, which corresponds to a composition of projections $\pi_n\circ\ldots\circ\pi_{k+1}=:\pi_{(n,k)}$. %This also allows us to see $W_C$ as the inverse limit of groups $G_n$ %We will be using quotients of this type and extensions of the presentations $G_n\to G_{n+1}$ described above to recursively define the presentation $\langle S \mid \mathcal U \rangle$. 

\smallskip

Now let $\alpha=\{[l_t,r_t]\}_{t=0}^\infty$ be a recursively enumerable sequence of segments from Proposition~\ref{prop-shaded-set}. We will use $\alpha$ to build the required presentation $\langle S \mid \mathcal U \rangle$ as a certain limit
%\footnote{As we consider both embeddings and quotients in the construction, the correct limit is inverse, but to make it work, one should extend generators of each presentation $P_k$ to a countable set $S=\{s_1,s_2\ldots \}$.} 
of an infinite sequence of finite presentations $\{P\}_{i=0}^{\infty}$. Each presentation $P_k$ is written in formal generators $\{s_1,s_2,\ldots,s_{f(c(k))}\}$, where $c(k)=\max\{r_i \mid i\leq k\}.$ The presentation $P_k$ presents a wreath product $G_{w(k)}$ for some function $w(k): \mathbb Z_{\geq 0} \to \mathbb Z_{\geq 0}$, which follows to be recursive from our construction, and subject to the inequality $w(k)\leq c(k)$. The group $G_{w(k+1)}$ is obtained from $G_{w(k)}$ by the extension composition $\psi_{w(k)}\circ\ldots\circ\psi_{w(k+1)-1}$ when $w(k+1)\geq w(k)$ or by the projection composition $\pi_{w(k)}\circ\ldots\circ\pi_{w(k+1)+1}$ when $w(k+1)<w(k)$. Before we proceed to the details, we introduce a partition of natural numbers into segments: 
\[
\mathbb N = \bigcup_{i=1}^{\infty} S_i, \ \ \text{where } S_i=[f(i-1)+1,f(i)].
\]
We will use this partition to split generators of $W_C=\langle \mathcal S \mid \mathcal U \rangle$ so that the generator indices of each base group in the iterated wreath product $W_C=\text{Wr}_{i\in \mathbb Z_-}C$ fit into a single segment.

%is the following correspondence between the generators of $P_k$ that are equal to $1$ and the shaded set of sequence $\alpha$:
%\begin{equation}
%\{ i \mid s_i=1 \text{ in } P_k \} = \bigcup \bigl\{ [E(p)+1,E(p+1)] \mid  p\in\mathcal{I}(\{[l_t,r_t]\}_{t=0}^k)  \bigr\}.
%\end{equation}

\smallskip

We start with a trivial group $G_0$ with the empty presentation $P_0$. At each step $k\in \mathbb Z_{\geq 0}$ we consider the segment $[l_k,r_k]$ and similarly to the cases in \eqref{it-case} we define $P_{k+1}$ using $P_k$. 

The key property we preserve is the following: $P_k$ has $f(c(k))=f(\max \{r_i \mid i\leq k\})$ generators, and for each segment $S_i$ either $s_j=1$ in $P_k$ for all $j\in S_i$ when $i\in \mathcal I(\{[l_t,r_t]\}_{t=0}^k)$; or some initial generators $s_j$ with $j\in S_i$ are not equal to $1$ in $P_k$, when $i\notin \mathcal I(\{[l_t,r_t]\}_{t=0}^k)$. 

\emph{Case 1: $l_{k+1}>\max \{r_i \mid i<k+1\}=c(k)$.} In this case we extend the presentation $P_k=\langle s_1,\ldots,s_{c(k)} \mid Q_k \rangle$ by adding new generators and relations. As $P_k$ presents the group $G_{w(k)}$, we extend this group so that $P_{k+1}$ presents $G_{w(k+1)}$, where $w(k+1)=w(k)+(l_{k+1}-c(k)-1)$. To do this, for each $i\in \{1,\ldots,l_{k+1}-c(k)-1\}$ we use the extension map $\psi_{w(k)+i}$ considering new generators $s_t$, where $t\in S_{c(k)+i}$. The first $dq^{w(k)+i-1}$ elements in this segment will be used for indexing the generators of the base group $B_{w(k)+i}$, and map the generators $s_t$ with remaining indices in $S_{w(k)+i}$ to identity. Since $|S_{c(k)+i}|=f(c(k)+i)-f(c(k)+i-1)=dq^{c(k)+i-1}$, and $w(k)\leq c(k)$, this allocation is possible. We also add the relations for the corresponding extension of $G_{w(k)+i-1}$ to $G_{w(k)+i}$ in the standard way. Finally, for each $i\in [l_{k+1},r_{k+1}]$ we add the relations $\{ s_j=1 \mid j \in  S_{i} \}$ ensuring  $c(k+1)=r_{k+1}$. It is easy to see that the key property is preserved by our construction and that the presentation $P_{k+1}$ defines the group $G_{w(k+1)}$.

\emph{Case 2: $l_{k+1}\leq\max \{r_i \mid i<k+1\}$.} Here we only need to factor out some generators in the presentation $P_k$ so that the key property is preserved. To do this, we consider generators $s_j$ with $j\in \{S_i \mid i\in [l_{k+1},c(k)]\}$ and add to $P_{k+1}$ a relation $s_j=1$ for each of them. Since the generators of the top base groups are fit into segments, the described relations factor out some top base groups in $G_{w(k)}$ seen as the iterated wreath product. We have that $c(k+1)=c(k)$ and the key property is preserved by \eqref{it-case} and our construction. One finds the number $w(k+1)$ by finding the size of the set $[1,c(k+1)]\setminus \mathcal I(\{[l_t,r_t]\}_{t=0}^{k+1})$.

\smallskip
Now we can take the countable set of generators $S=\{s_1,s_2,\ldots \}$ and consider the presentation $\langle S \mid \mathcal U\rangle$, where $\mathcal U = \bigcup_{k=1}^{\infty} Q_k$. Since the set $\mathbb N\setminus \mathcal{I}(\alpha)$ is infinite, there will be infinitely many $k\in \mathbb N$ such that the presentation $P_k=\langle s_1,s_2,\ldots,s_{f(c(k))} \mid Q_k \rangle$ remains stable in the union, so we have $G_k<\langle S \mid \mathcal U \rangle$. By considering an increasing subsequence of such indices $k(t)$ we obtain $\langle S \mid \mathcal U \rangle = \lim_{t\to \infty} \langle S\mid Q_{k(t)} \rangle \cong W_C$. 

It remains to notice that the key property of the construction implies that the word problem of $\langle S \mid \mathcal U\rangle$ is reducible to the presence problem for the set $\mathcal{I}(\alpha)$. Indeed, if $i$ is the first element in the segment $S_j$, then $s_i=1$ is equivalent to checking whether $j\in\mathcal I(\alpha)$. Since the latter is undecidable by the choice of the sequence $\alpha$, the result follows.
\end{proof}

\end{document}
